%% LyX 2.4.4 created this file.  For more info, see https://www.lyx.org/.
%% Do not edit unless you really know what you are doing.
\documentclass[english]{article}
\usepackage[T1]{fontenc}
\usepackage[latin9]{inputenc}
\usepackage{units}
\usepackage{enumitem}
\usepackage{amsmath}
\usepackage{amsthm}
\usepackage{amssymb}
\usepackage{geometry}
\geometry{verbose,tmargin=1in,bmargin=1in,lmargin=1in,rmargin=1in}

\makeatletter
%%%%%%%%%%%%%%%%%%%%%%%%%%%%%% Textclass specific LaTeX commands.
\numberwithin{equation}{section}
\numberwithin{figure}{section}
\newlength{\lyxlabelwidth}      % auxiliary length 

%%%%%%%%%%%%%%%%%%%%%%%%%%%%%% User specified LaTeX commands.
\usepackage{fancyhdr}
\usepackage{lastpage}
\pagestyle{fancy}
\usepackage{enumitem}
\usepackage{ifthen}
\everymath=\expandafter{\the\everymath\displaystyle}
%\DeclareMathSizes{10}{10}{10}{10}
\usepackage{multicol}

\usepackage{tikz}
\usepackage{tikz-3dplot}
\usetikzlibrary{calc,arrows}

\usetikzlibrary{patterns}
\usetikzlibrary{plotmarks}
\usepackage{pgfplots}
\pgfplotsset{compat=newest}

\usepackage{calc}
\usepackage[nomessages]{fp}

\usepackage{datetime}
% Added by lyx2lyx
\setlength{\parskip}{\smallskipamount}
\setlength{\parindent}{0pt}

\makeatother

\usepackage{babel}
\begin{document}
\renewcommand{\author}{Dr.\,James Ashe}

\newcommand{\classNum}{336}

\newcommand{\classSec}{A}

\newcommand{\examType}{Exam 2}

\newcommand{\nameType}{Name:}

\renewcommand{\date}{\formatdate{25}{10}{2013}}

%%First page's header%%%%%%%%%%%%%%%%%%%%%%
\fancypagestyle{firststyle} %%header settings for first pagr
{
	\fancyhead[L]{MTH \classNum{}(\classSec{}) \author{}\\
	\textbf{\examType{}} ~\date{}}
	\fancyhead[C]{\textbf{\nameType}}
	\fancyhead[R]{}
	\setlength{\headsep}{14pt}
}
%%%%%%%%%%%%%%%%%%%%%%%%%%%%%%%%%%%%%%%%%%

%%Next page's header%%%%%%%%%%%%%%%%%%%%%%%
\setlist{nolistsep}
\lhead{MTH \classNum{}(\classSec{})} 
\chead{\textbf{\nameType}} 
\cfoot{\thepage{}~of~\pageref{LastPage}} 
\setlength{\headsep}{5pt} 
%%%%%%%%%%%%%%%%%%%%%%%%%%%%%%%%%%%%%%%%%%%

%%Various settings; see the preamble for other settings%%%%%
%\setenumerate[0]{leftmargin=12pt,itemindent=0pt,labelwidth=0pt}
\setlist{topsep=.5pt}
\renewcommand{\labelenumi}{\textbf{\arabic{enumi}.}}
\renewcommand{\labelenumii}{\textbf{\alph{enumii}.}}
\thispagestyle{firststyle}
%%%%%%%%%%%%%%%%%%%%%%%%%%%%%%%%%%%%%%%%%%

\newcommand{\xmax}{0}
\newcommand{\xmin}{-\xmax}
\newcommand{\xstep}{0}
\newcommand{\ymax}{0}
\newcommand{\ymin}{-\ymax}
\newcommand{\ystep}{0} 
\newcommand{\dspace}{\vspace{1cm}}
\newcommand{\xa}{0}
\newcommand{\xb}{0}
\newcommand{\ya}{1}
\newcommand{\yb}{1}

\newcommand{\xspace}{\vspace{.75cm}}

\newcommand{\vectors}[2]{\textbf{#1}_{1},\textbf{#1}_{2},\dots,\textbf{#1}_{#2}}
\newcommand{\vectorsN}[1]{\textbf{#1}_{1},\textbf{#1}_{2},\dots,\textbf{#1}_{n}}
\newcommand{\R}[1]{\mathbb{R}^{#1}}
\begin{enumerate}[resume]
\item Let $A=\begin{bmatrix}\begin{array}{rr}
1 & 2\\
3 & 4
\end{array}\end{bmatrix}$, $B=\begin{bmatrix}\begin{array}{rr}
-1 & 1\\
0 & 1
\end{array}\end{bmatrix}$, and $C=\begin{bmatrix}\begin{array}{rrr}
1 & 2 & 2\\
0 & 2 & 4
\end{array}\end{bmatrix}$. Calculate each of the following expressions or explain why it is
not defined.

\begin{enumerate}[resume]
\item $AB$\xspace \vfill
\item $CB$\xspace
\end{enumerate}
\item Suppose that $T\colon\R{3}\rightarrow\R{2}$ is defined by $T\left(\begin{bmatrix}\begin{array}{r}
x_{1}\\
x_{2}\\
x_{3}
\end{array}\end{bmatrix}\right)=\begin{bmatrix}\begin{array}{r}
2x_{1}\\
x_{2}+x_{3}
\end{array}\end{bmatrix}$.

\begin{enumerate}
\item Find the standard matrix $A$ for $T$.\xspace \xspace \vfill
\end{enumerate}
\begin{enumerate}[resume]
\item Is $T$ is a linear transformation.\xspace
\end{enumerate}
\item Let $\mathbf{x}\in\R{n}$ and $A$ be an $m\times n$ matrix. The
function $f(\mathbf{x})=A\mathbf{x}$ is a linear map from \underline{\hspace{.5cm}}
to \underline{\hspace{.5cm}}.\xspace
\item Find $A^{-1}$ or show that it does not exist.

\begin{enumerate}
\item $A=\begin{bmatrix}1 & 1 & 0\\
0 & 1 & 1\\
0 & 0 & 1
\end{bmatrix}$\xspace \xspace \vfill
\end{enumerate}
\begin{enumerate}[resume]
\item $A=\begin{bmatrix}2 & 2 & 0\\
0 & 2 & 2\\
0 & 0 & 0
\end{bmatrix}$\xspace
\end{enumerate}
\item Let $A=\begin{bmatrix}\begin{array}{rr}
\nicefrac{1}{2} & \nicefrac{1}{4}\\
0 & \nicefrac{1}{2}
\end{array}\end{bmatrix}$, $A^{-1}=\begin{bmatrix}\begin{array}{rr}
2 & -1\\
0 & 2
\end{array}\end{bmatrix}$ and $\mathbf{b}=\begin{bmatrix}\begin{array}{r}
1\\
1
\end{array}\end{bmatrix}$

\begin{enumerate}
\item Use $A^{-1}$ to solve $A\mathbf{x}=\mathbf{b}$ for \textbf{$\mathbf{x}$}.\xspace \vfill
\end{enumerate}
\begin{enumerate}[resume]
\item Use your answer from part \textbf{(a)} to express $\mathbf{b}$ as
a linear combination of $A$.\xspace
\item Find $(A^{-1})^{-1}$ or show that it does not exist.\xspace
\end{enumerate}
\item Let $\mathbf{v}_{1}=(1,2)$, $\mathbf{v}_{2}=(0,1)$, and $\mathbf{v}_{2}=(1,0)$
be three vectors in $\R{2}$.

\begin{enumerate}
\item Show that $\mathbf{v}_{1}=(1,2)$ and $\mathbf{v}_{2}=(0,1)$ gives
a basis for $\R{2}$.\xspace \vfill
\end{enumerate}
\begin{enumerate}[resume]
\item Find another basis for $\R{2\}}$.\xspace
\item Is $\mathbf{v}_{1}=(1,2)$, $\mathbf{v}_{2}=(0,1)$, and $\mathbf{v}_{2}=(1,0)$
linearly independent? Explain.\xspace
\item Is $\mathbf{v}_{1}=(1,2)$, $\mathbf{v}_{2}=(0,1)$, and $\mathbf{v}_{2}=(1,0)$
a basis for $\R{2}$? Explain.\xspace
\item Does $\mathbf{v}_{1}=(1,2)$, $\mathbf{v}_{2}=(0,1)$, and $\mathbf{v}_{2}=(1,0)$
span $\R{2}$? Explain. \xspace
\end{enumerate}
\item Let $\mathbf{v}_{1}=(2,4,8)$, $\mathbf{v}_{2}=(-2,-4-8)$, $\mathbf{v}_{2}=(1,2,4)$
be vectors in $\R{3}$. 

\begin{enumerate}
\item Is $\mbox{Span}(\mathbf{v}_{1},\mathbf{v}_{2},\mathbf{v}_{3})$ a
subspace of $\R{3\}}$.\xspace
\end{enumerate}
\begin{enumerate}[resume]
\item Is $\{\mathbf{v}_{1},\mathbf{v}_{2},\mathbf{v}_{3}\}$ linearly independent?\xspace \vfill
\item Does $\{\mathbf{v}_{1},\mathbf{v}_{2},\mathbf{v}_{3}\}$ span $\R{3}$?,
i.e., does $\mbox{Span}(\mathbf{v}_{1},\mathbf{v}_{2},\mathbf{v}_{3})=\R{3}$?\xspace
\item Find the dimension of $\mbox{Span}(\mathbf{v}_{1},\mathbf{v}_{2},\mathbf{v}_{3})$.\xspace
\item Find a basis for $\mbox{Span}(\mathbf{v}_{1},\mathbf{v}_{2},\mathbf{v}_{3})$.
\xspace
\end{enumerate}
\item Let $V=\{\vectors{v}{n+1}\}$ be a set of $n+1$ vectors in $\mathbb{R}^{n}$.

\begin{enumerate}
\item Is $V$ linearly independent or linearly dependent?\xspace
\end{enumerate}
\begin{enumerate}[resume]
\item Is $V$ a basis for $\mathbb{R}^{n}$? Explain. \xspace
\item Is $\mbox{span}V$ a subspace of $\mathbb{R}^{n}$. Explain.\xspace
\item Is $V$ necessarily a subspace of $\mathbb{R}^{n}$, i.e., is a set
of $n+1$ vectors in $\mathbb{R}^{n}$ always a subspace of $\mathbb{R}^{n}$?
Explain.\xspace
\end{enumerate}
\item \textbf{Bonus.}Let $\mathbf{v}_{1},\mathbf{v}_{2},\dots,\mathbf{v}_{k}\in\mathbb{R}^{n}$
and let $\mathbf{v}\in\mathbb{R}^{n}$. Prove that $\mathbf{v}\in\mbox{Span}\left(\mathbf{v}_{1},\mathbf{v}_{2},\dots,\mathbf{v}_{k}\right)$
$\iff$ $\mbox{Span}\left(\mathbf{v}_{1},\mathbf{v}_{2},\dots,\mathbf{v}_{k}\right)=\mbox{Span}\left(\mathbf{v}_{1},\mathbf{v}_{2},\dots,\mathbf{v}_{k},\mathbf{v}\right)$.
\end{enumerate}

\end{document}
